% Generated draft: requires mathematical and source review.
\documentclass[11pt]{article}
\usepackage[margin=25mm]{geometry}
\usepackage{fontspec}
\setmainfont{Latin Modern Roman}
\usepackage{amsmath,amssymb,mathtools,unicode-math}
\setmathfont{Latin Modern Math}
\usepackage{enumitem,needspace,xurl,hyperref}
\setlength{\parindent}{0pt}
\setlength{\parskip}{5pt}
\setlength{\emergencystretch}{3em}
\newcommand{\N}{\mathbb{N}}
\newcommand{\Z}{\mathbb{Z}}
\newcommand{\Q}{\mathbb{Q}}
\newcommand{\R}{\mathbb{R}}
\newcommand{\C}{\mathbb{C}}
\newcommand{\F}{\mathbb{F}}
\newcommand{\A}{\mathbb{A}}
\newcommand{\PP}{\mathbb{P}}
\newcommand{\E}{\mathbb{E}}
\newcommand{\eps}{\varepsilon}
\newcommand{\dd}{\,\mathrm{d}}
\begin{document}
\section*{UnsolvedMath ID: 800036}
\textbf{Source problem code:} catalogue entry
\subsection{Definitions and mathematical statement}
Let \(S\) be a semigroup, that is, a nonempty set with an associative binary operation. A group \(G\) is periodic if every element has finite order:
\[
\forall g\in G\;\exists n\in\Z_{\geq 1}\quad g^n=1_G.
\]
For subsets \(A,B\subseteq G\), define their setwise product by
\[
AB=\{ab:a\in A,\ b\in B\}.
\]
Two conventions for the power semigroup are
\[
\mathcal P_{+}(G)=\{A\subseteq G:A\ne\varnothing\},
\qquad
\mathcal P_{\mathrm{all}}(G)=\{A:A\subseteq G\},
\]
each equipped with setwise multiplication. Under the second convention, the empty set is an absorbing element. The supplied record does not specify which convention is intended.

For a fixed convention \(c\in\{+,\mathrm{all}\}\), the embeddability property in question is the existence of a periodic group \(G\) and an injective map \(\phi:S\to\mathcal P_c(G)\) such that
\[
\phi(st)=\phi(s)\phi(t)\qquad\text{for all }s,t\in S.
\]
The ambient group may depend on \(S\). This is an embedding of semigroups; no additional identity-preservation requirement is stated in the source.

\textbf{Problem (S. G. Bershadsky and S. I. Kublanovsky).} Find general necessary and/or sufficient conditions on a semigroup \(S\) for this embeddability property. The intended convention concerning the empty subset remains unresolved, so the two definitions above specify alternative readings of the source rather than an assertion that the readings are equivalent. The supplied record reports that some partial conditions were obtained in references labelled [Donald] and [Bersh+Kubl], but does not state those conditions or provide their bibliographic details.
\subsection{Short English statement}
Find general conditions ensuring, or required for, a semigroup to embed into the semigroup of subsets of some group whose elements all have finite order, with multiplication given by setwise products. The source leaves unclear whether the empty subset is included. It reports partial conditions in [Donald] and [Bersh+Kubl].
\subsection{Sources}
\begin{itemize}
\item Local UnsolvedMath registry, permanent ID \(800036\), exact source problem code \texttt{catalogue entry}, title ``The Lyapin's Notebook --- Problem 38'', in the supplied AMR Open Problem Lists record. The registry labels the problem open; this is a supplied status, not an independently verified assessment.
\item Supplied source list: Bershadsky, Kublanovsky, Mashevitzky, \emph{The Lyapin's Notebook: Unsolved Problems in Semigroup Theory} (2026), Problem 38. Supplied URL: \url{https://arxiv.org/abs/2604.04763}. The record reports access on 2026-10-16 and extraction by \texttt{arxiv\_latex}; these are supplied provenance statements, not a claim of independent source consultation.
\item The problem statement attributes this problem specifically to S. G. Bershadsky and S. I. Kublanovsky. The registry's broader proposer field lists Bershadsky, Kublanovsky, Mashevitzky and records the year 2026.
\item [Donald] and [Bersh+Kubl]: source reference labels cited for partial conditions. No full bibliographic entries, links, or statements of the partial conditions are supplied.
\end{itemize}
\end{document}
